\subsection{\textit{Modeling} dan Validasi}
\label{subsec:lr-validasi}

\subsubsection{Skema Validasi}
\todo{Skema validasi untuk LR (\textit{stratified k-fold} atau \textit{train-test split}), nilai $k$/rasio, dan \textit{seed}. Jika sama dengan kerangka umum (Subbab~\ref{sec:kerangka-umum}), cukup rujuk dan jelaskan perbedaannya saja, misalnya \textit{subsampling} karena waktu komputasi.}

\subsubsection{\textit{Hyperparameter Tuning}}
\todo{Ruang pencarian (jenis penalti, $\lambda$, $\rho$ (\textit{Elastic Net}), \textit{learning rate}, \textit{batch size}, jumlah \textit{epoch}), metode pencarian, dan hasilnya. Sertakan tabel atau grafik skor CV terhadap hiperparameter.}

\begin{table}[H]
	\centering
	\caption{Hasil \textit{tuning} hiperparameter LR}
	\label{tab:lr-tuning}
	\begin{tabular}{l l c}
		\toprule
		\textbf{Hiperparameter} & \textbf{Ruang pencarian} & \textbf{Terbaik} \\
		\midrule
		\todo{} & & \\
		\bottomrule
	\end{tabular}
\end{table}

\subsubsection{Model Final}
\todo{Konfigurasi final LR, skor CV (rata-rata $\pm$ simpangan baku), dan lokasi model tersimpan di \texttt{output/}.}
